\chapter{Prolactin}

\section*{What Is This Test?}
Prolactin is a 199-amino-acid single-chain polypeptide hormone produced by lactotrope cells of the anterior pituitary gland, which constitute approximately 15--25\% of anterior pituitary cells. Prolactin's primary physiological role is the initiation and maintenance of lactation (milk production) in the postpartum period. It also plays a role in reproductive function by inhibiting gonadotropin-releasing hormone (GnRH) secretion, thereby suppressing ovulation during breastfeeding (lactational amenorrhea). Prolactin secretion is unique among anterior pituitary hormones because it is under tonic inhibitory control by dopamine (prolactin-inhibiting factor) produced by the hypothalamic tuberoinfundibular dopaminergic neurons and delivered to the anterior pituitary via the hypothalamic-hypophyseal portal system. Dopamine binds to D2 receptors on lactotrope cells, inhibiting prolactin synthesis and secretion. Thyrotropin-releasing hormone (TRH) and vasoactive intestinal peptide (VIP) stimulate prolactin release. Prolactin exhibits a circadian rhythm with peak levels during sleep (highest at 4--7 AM) and a nadir in the late morning, and it is also stimulated by stress, exercise, breast stimulation, and meals. Prolactin circulates in multiple molecular forms: monomeric prolactin (23 kDa, 80--90\% of circulating prolactin, biologically active), big prolactin (50--60 kDa, dimeric), and macroprolactin (>150 kDa, prolactin complexed with IgG autoantibodies, biologically inactive but detected by immunoassays). The prolactin test is primarily used to investigate galactorrhea, infertility, menstrual disorders, and suspected prolactinoma.

\section*{Alternative Names}
PRL; Serum Prolactin; Lactogenic Hormone; Lactotropin; Mammotropin; Macroprolactin; Prolactin, Monomeric; Prolactin Fractionation by PEG Precipitation

\section*{Principle}
Prolactin is measured using a two-site sandwich immunometric (chemiluminescent) immunoassay on automated analyzers. Two monoclonal antibodies directed against different epitopes on the prolactin molecule are used: a capture antibody bound to a solid phase (magnetic particles or microtiter wells), and a detection antibody labeled with a chemiluminescent tag (ruthenium on Roche ECLIA, acridinium ester on Abbott/Siemens, alkaline phosphatase on Beckman Coulter). After incubation and washing, the signal is proportional to prolactin concentration. The assay detects all molecular forms of prolactin (monomeric, big, and macroprolactin) but is standardized against monomeric prolactin (WHO 3rd IS 84/500). Functional sensitivity is typically <0.5 ng/mL \cite{tietz2012textbook}. A major clinical challenge is macroprolactin, a complex of monomeric prolactin with IgG autoantibodies that is biologically inactive but detected as "prolactin" by all commercial immunoassays. Macroprolactin accounts for 15--25\% of all cases of hyperprolactinemia and the vast majority of asymptomatic cases. Macroprolactin can be identified by polyethylene glycol (PEG) precipitation: PEG precipitates high-molecular-weight complexes (macroprolactin), and the prolactin recovered in the supernatant (monomeric prolactin) is measured. If post-PEG recovery is <40--50\% of the total prolactin, macroprolactin is the main circulating form, and the patient is likely euvolemic (does not have true hyperprolactinemia). Alternatively, gel filtration chromatography can separate prolactin isoforms, but this is not available in routine clinical laboratories.

\section*{History}
Prolactin was the last anterior pituitary hormone to be definitively identified as a distinct hormone from growth hormone. Its existence was suspected from the 1930s, but it was not until 1970 that Henry Friesen and colleagues definitively isolated and purified human prolactin, demonstrating it was distinct from GH \cite{friesen1970biosynthesis}. Friesen also developed the first radioimmunoassay for prolactin in 1971. The dopamine-prolactin axis was elucidated in the 1970s, demonstrating that dopamine is the primary prolactin-inhibiting factor. The clinical syndrome of hyperprolactinemia (amenorrhea-galactorrhea syndrome) was described by Forbes, Henneman, and colleagues in the 1950s, and the association with pituitary tumors was established in the 1970s. Prolactin-secreting adenomas (prolactinomas) were recognized as the most common type of functional pituitary adenoma. Dopamine agonists for treating prolactinomas were introduced: bromocriptine in the 1970s and cabergoline in the 1990s (cabergoline has a longer half-life and better tolerability, making it first-line therapy). Macroprolactin was described in the 1980s and its prevalence (15--25\% of hyperprolactinemic patients) and clinical insignificance were established in the 1990s--2000s. The 2011 Endocrine Society clinical practice guideline on hyperprolactinemia diagnosis and treatment standardized the approach to prolactin measurement, including the recommendation for macroprolactin screening \cite{melmed2011diagnosis}.

\section*{Why This Test Is Ordered}
\begin{itemize}
  \item Investigation of galactorrhea (milky nipple discharge not associated with pregnancy or lactation) in women --- present in 30--80\% of hyperprolactinemic women, depending on estrogen status
  \item Evaluation of menstrual disturbances: amenorrhea (secondary or primary), oligomenorrhea, or infertility in women --- hyperprolactinemia suppresses GnRH, leading to anovulation and estrogen deficiency
  \item Investigation of male hypogonadism: decreased libido, erectile dysfunction, infertility, gynecomastia, and rarely galactorrhea --- hyperprolactinemia suppresses GnRH, reducing LH and testosterone
  \item Diagnosis and monitoring of prolactin-secreting pituitary adenomas (prolactinomas) --- serum prolactin levels correlate with tumor size: microprolactinomas (<1 cm) typically have prolactin 50--200 ng/mL; macroprolactinomas (>1 cm) typically have prolactin >200--500 ng/mL, often >1000 ng/mL
  \item Investigation of the "stalk effect" --- non-functioning pituitary adenomas or other sellar/suprasellar masses compressing the pituitary stalk disrupt dopaminergic inhibition of lactotropes, causing mild-to-moderate hyperprolactinemia (typically <100--150 ng/mL, distinguishing from prolactinoma where prolactin is usually >200 ng/mL for tumors >1 cm and >500 ng/mL for tumors >2 cm)
  \item Screening for macroprolactin in patients with asymptomatic hyperprolactinemia --- PEG precipitation or gel filtration chromatography \cite{gibney2005impact}
  \item Monitoring dopamine agonist therapy (cabergoline, bromocriptine) in treated prolactinoma patients --- target is normoprolactinemia and tumor shrinkage
  \item Evaluation of headache, visual field defects, or hypopituitarism associated with pituitary/sellar mass --- prolactin is part of the pituitary hormone panel
  \item Investigation of drug-induced hyperprolactinemia from antipsychotics, antidepressants, antiemetics, or other dopamine antagonist medications
\end{itemize}

\section*{Primary vs Secondary Test}
Prolactin is a primary screening test for hyperprolactinemia when patients present with galactorrhea, amenorrhea/oligomenorrhea, infertility, or male hypogonadism. It is part of the basic pituitary function panel in the evaluation of sellar masses and suspected hypopituitarism. Because stress and venipuncture can cause mild, transient prolactin elevation, a mildly elevated result should be confirmed by repeat testing under non-stressed, fasting conditions before proceeding to extensive evaluation.

\section*{How This Test Is Performed}
A healthcare professional draws a blood sample from a vein into a serum separator tube (gold-top) or plain red-top tube. The sample should be drawn in the morning (before 10 AM), preferably 3--4 hours after waking, to avoid the sleep-related prolactin peak. The patient should be seated and relaxed for 15--30 minutes before venipuncture, because stress and venipuncture itself can transiently elevate prolactin. The specimen does not require special pre-analytical handling (no need for ice or freezing). Prolactin is stable at room temperature for 8 hours, at 2--8 degrees C for 7 days, and long-term at -20 degrees C. For macroprolactin screening, an aliquot of the serum sample is treated with 25\% PEG solution to precipitate high-molecular-weight complexes; the supernatant is then assayed for prolactin. Post-PEG recovery <40--50\% indicates macroprolactin predominance. Macroprolactin screening should be performed in all patients with hyperprolactinemia, particularly if asymptomatic.

\section*{What Preparation Is Needed}
The patient should fast for 8--10 hours before testing because meals stimulate prolactin release. The blood draw should occur in the morning (8--10 AM) to avoid the nocturnal prolactin surge. The patient should avoid breast stimulation, chest wall trauma, and sexual intercourse for 24 hours before testing because nipple stimulation triggers prolactin release through a neuroendocrine reflex. Strenuous exercise should be avoided for 12 hours because it elevates prolactin. The patient should be awake for at least 3--4 hours before blood draw. Stress and anxiety during venipuncture can elevate prolactin --- if the initial result is mildly elevated, a repeat sample should be drawn through an indwelling catheter after 30 minutes of rest (the "rested prolactin"). Several medications that elevate prolactin should ideally be discontinued for 2--4 weeks before testing if clinically safe: antipsychotics (especially typical agents like haloperidol and atypical agents like risperidone --- all dopamine D2 antagonists); metoclopramide and domperidone (D2 antagonists used as antiemetics/prokinetics) --- markedly elevate prolactin 5--10 fold within hours of a single dose; tricyclic antidepressants and SSRIs (fluoxetine, paroxetine, sertraline) mildly elevate prolactin; verapamil (calcium channel blocker) mildly elevates prolactin; H2-receptor antagonists (cimetidine) mildly elevate prolactin; opioids; estrogen therapy (oral contraceptives, HRT) mildly elevate prolactin (by 20--30\%). High-dose biotin (>5 mg/day) must be discontinued 48 hours before testing.

\section*{Contraindications and Precautions}
Factors causing false-positive elevations (transient, physiologic): stress of venipuncture (very common, can increase prolactin 2--3 fold --- a repeat rested sample is essential before extensive evaluation); recent exercise, sexual intercourse, or breast stimulation (<24 hours); recent meal (mild elevation); sleep deprivation or recent awakening (prolactin peaks during REM sleep); pregnancy and lactation (prolactin is physiologically elevated 5--10 fold during pregnancy and remains elevated during breastfeeding, declining 4--6 weeks after cessation); chest wall trauma, herpes zoster (thoracic dermatomes), thoracotomy, or nipple piercings can stimulate prolactin release through the neuroendocrine reflex. Drugs causing hyperprolactinemia include (most clinically important): antipsychotics (all D2 antagonists, especially risperidone, haloperidol, paliperidone; the most common drug cause of clinically significant hyperprolactinemia --- prolactin may be 100--300 ng/mL); metoclopramide, domperidone (D2 antagonists, elevate prolactin even after a single dose); antidepressants (mild prolactin elevation, typically <50 ng/mL); verapamil (mild, <50 ng/mL); opiates; cocaine; alpha-methyldopa. \cite{molitch2008drugs} Pathologic false-positives: primary hypothyroidism (elevated TRH stimulates prolactin, prolactin typically 30--100 ng/mL and normalizes with levothyroxine); chronic kidney disease (decreased renal clearance of prolactin, prolactin 50--200 ng/mL in CKD stages 4--5); cirrhosis (alterations in dopaminergic tone); polycystic ovary syndrome (PCOS) (mild elevation, prolactin <50 ng/mL, related to relative estrogen excess without progesterone opposition); epilepsy (ictal and post-ictal prolactin elevation, used to differentiate epileptic seizures from psychogenic non-epileptic seizures). Macroprolactin: biologically inactive but detected by all immunoassays, accounts for 15--25\% of all hyperprolactinemia cases, typically completely asymptomatic (no galactorrhea, regular menses). False-negative from hook effect (prozone phenomenon): in giant macroprolactinomas with prolactin >5000--10,000 ng/mL, extremely high prolactin concentrations saturate both capture and detection antibodies in the sandwich immunoassay, preventing sandwich formation and producing falsely low/normal results. This is avoided by diluting the sample 1:100 and re-assaying. The hook effect should be excluded in any patient with a large (>3 cm) pituitary mass and normal or modestly elevated prolactin.

\section*{Risks}
Minimal risk. Slight pain, bruising, or hematoma at venipuncture site. Rarely: vasovagal syncope. No radiation exposure. Macroprolactin screening by PEG precipitation: no additional risk (performed on serum aliquot).

\section*{What Happens After the Test}
Results are available within 1--4 hours on automated analyzers. If prolactin is mildly elevated (<100 ng/mL), the first step is to repeat the measurement in a resting, fasting state, after excluding interfering medications \cite{vilar2008diagnosis}. If confirmed hyperprolactinemia, TSH, creatinine, and pregnancy test should be obtained to exclude secondary causes. Macroprolactin should be excluded by PEG precipitation, especially if the patient is asymptomatic. If macroprolactin is confirmed and the patient is asymptomatic, no further evaluation or treatment is needed. If true monomeric hyperprolactinemia is confirmed and secondary causes (medications, hypothyroidism, CKD, pregnancy) are excluded, a pituitary MRI with gadolinium contrast should be performed to evaluate for a prolactinoma or other sellar mass. Prolactin levels >200 ng/mL are almost always due to a prolactinoma; levels <100 ng/mL may be due to stalk compression from a non-functioning mass, drug effect, or microprolactinoma. For confirmed prolactinoma, treatment with dopamine agonist (cabergoline 0.25--0.5 mg twice weekly, titrated up based on prolactin levels and tumor size) is first-line therapy, normalizing prolactin in 80--90\% and shrinking tumor in 70--80\% \cite{klibanski2010prolactinomas}.

\section*{Understanding Results}

\subsection*{Below Normal}
\begin{itemize}
  \item \textbf{Low prolactin (<3--5 ng/mL in non-pregnant, non-lactating women; <2 ng/mL in men):} Hypoprolactinemia. Clinically significant only in the postpartum period, when low prolactin causes failure of lactation (Sheehan syndrome --- postpartum pituitary necrosis from obstetric hemorrhage). Other causes: panhypopituitarism (destruction of lactotrope cells), dopamine agonist therapy (cabergoline, bromocriptine --- pharmacologically suppresses prolactin), and isolated prolactin deficiency (rare, associated with POU1F1/Pit-1 mutations).
  \item \textbf{Suppressed prolactin in Sheehan syndrome:} Postpartum hemorrhage causing pituitary infarction/necrosis. Failure of postpartum lactation is the earliest and most characteristic feature. Other pituitary hormones may be affected (ACTH > GH > gonadotropins > TSH, in order of vulnerability).
  \item \textbf{Post-surgical remission of prolactinoma:} Surgical resection (typically transsphenoidal, reserved for dopamine agonist-resistant or intolerant macroprolactinomas, or microprolactinomas in patients wishing to avoid long-term dopamine agonist therapy) should normalize prolactin. Prolactin <5 ng/mL 1 week post-operative indicates likely surgical cure.
  \item \textbf{Cabergoline/bromocriptine-induced suppression:} Expected effect of dopamine agonist therapy. Target is prolactin in the normal range. Over-suppression is not harmful; the dose can be reduced if prolactin is undetectable and the patient is asymptomatic.
\end{itemize}

\subsection*{Above Normal}
\begin{itemize}
  \item \textbf{Prolactinoma (prolactin-secreting pituitary adenoma):} The most common cause of pathologic hyperprolactinemia and the most common functional pituitary adenoma (40--45\% of all pituitary adenomas). Prolactin correlates with tumor size: microprolactinoma (<1 cm) --- prolactin 50--200 ng/mL (often <200 ng/mL); macroprolactinoma (>1 cm) --- prolactin >200--500 ng/mL (often >1000 ng/mL for tumors >2 cm). Symptoms in premenopausal women: galactorrhea (30--80\%), oligomenorrhea/amenorrhea, infertility, decreased libido. In men and postmenopausal women: mass effects (headache, visual field defects --- bitemporal hemianopsia from optic chiasm compression), hypogonadism, decreased libido, erectile dysfunction, and rarely galactorrhea (men lack sufficient estrogen priming). Prolactin levels >5000 ng/mL suggest a giant invasive macroprolactinoma. Treatment: dopamine agonists (cabergoline is first-line) \cite{wong2015update2}.
  \item \textbf{Medication-induced hyperprolactinemia:} The most common overall cause of elevated prolactin in clinical practice. Antipsychotics (all D2 antagonists): risperidone, paliperidone, haloperidol cause the most marked elevations (100--300 ng/mL); atypical antipsychotics with lower D2 affinity (olanzapine, quetiapine, clozapine, aripiprazole) cause minimal prolactin elevation and may be used as alternatives. Metoclopramide and domperidone: potent D2 antagonists, prolactin may rise within hours of a single dose to 100--200 ng/mL. Antidepressants (SSRIs, TCAs): mild elevation, typically <50 ng/mL, rarely >100 ng/mL. Verapamil: mild elevation, <50 ng/mL. If the medication cannot be discontinued, a dopamine agonist can be added (under endocrinology supervision, particularly if a prolactinoma also needs to be excluded).
  \item \textbf{Stalk effect (pituitary stalk compression):} Non-functioning pituitary adenomas or other sellar/suprasellar masses (craniopharyngioma, Rathke cleft cyst, meningioma, metastases) compress the pituitary stalk, disrupting the delivery of dopamine from the hypothalamus to the lactotropes. Loss of tonic inhibition causes mild-to-moderate hyperprolactinemia, typically <100--150 ng/mL (rarely >200 ng/mL). This is a critical distinction from a true prolactinoma: a large sellar mass (>2 cm) with prolactin <200 ng/mL suggests stalk effect from a non-functioning adenoma, NOT a prolactinoma (which would produce prolactin >500--1000 ng/mL for a tumor that size). An important clinical pitfall is mistaking stalk-effect hyperprolactinemia for a prolactinoma and treating with a dopamine agonist, which would shrink the normal lactotrope tissue but not the non-functioning tumor.
  \item \textbf{Primary hypothyroidism:} Elevated TRH stimulates lactotropes to secrete prolactin. TSH is elevated, free T4 is low. Prolactin elevation is mild-to-moderate (typically 30--100 ng/mL). Prolactin normalizes with levothyroxine therapy, and no pituitary imaging or dopamine agonist is needed.
  \item \textbf{Macroprolactinemia:} Macroprolactin (prolactin-IgG complex) is biologically inactive but detected by all commercial prolactin immunoassays. Affects 15--25\% of hyperprolactinemic individuals. Patients are typically asymptomatic (normal menstruation, no galactorrhea, normal fertility). Post-PEG precipitation recovery <40--50\% confirms macroprolactinemia. No treatment or further investigation is needed. Recognized as a common cause of unnecessary pituitary imaging and treatment.
  \item \textbf{Chronic kidney disease (CKD stages 4--5):} Decreased renal clearance of prolactin (kidneys clear approximately 25\% of circulating prolactin) and altered hypothalamic dopaminergic tone. Prolactin typically 50--200 ng/mL. Men may develop gynecomastia and hypogonadism; women may develop amenorrhea.
  \item \textbf{Lactation and pregnancy:} Physiologic hyperprolactinemia. Prolactin rises throughout pregnancy (50--300 ng/mL at term) and remains elevated during breastfeeding (basal levels 50--200 ng/mL, with surges to 200--500 ng/mL during suckling). Return to normal 4--6 weeks after cessation of breastfeeding.
  \item \textbf{Stress-induced hyperprolactinemia:} Stress (venipuncture, anxiety, illness, surgery, hypoglycemia) can elevate prolactin 2--3 fold through suppression of hypothalamic dopamine. This is the most common cause of mild, transient hyperprolactinemia. A repeat rested, fasting morning prolactin normalizes in >50\% of cases.
  \item \textbf{Chest wall/thoracic lesions:} Herpes zoster involving thoracic dermatomes, thoracotomy scars, burns, trauma, or nipple piercings can stimulate prolactin through the neuroendocrine suckling reflex arc (afferent impulses from T4--T6 dermatomes).
  \item \textbf{Seizures:} Generalized tonic-clonic seizures (not simple/partial seizures) cause a marked prolactin surge within 10--20 minutes post-ictally. Prolactin measured 10--20 minutes after a suspected seizure >3 times baseline is approximately 90\% specific for epileptic seizure vs. psychogenic non-epileptic seizure (PNES). Prolactin normalizes within 1--2 hours.
  \item \textbf{Hook effect (prozone phenomenon):} In giant macroprolactinomas with prolactin >5000--10,000 ng/mL, the sandwich immunoassay produces a falsely low/normal result because both antibodies are saturated by excess antigen, preventing sandwich formation. This is a dangerous diagnostic pitfall: a patient with a large sellar mass (>3 cm) and "normal" prolactin may be misdiagnosed with a non-functioning adenoma (which would be treated surgically) instead of a prolactinoma (which is treated medically with dopamine agonists). All patients with sellar masses >3 cm should have prolactin measured on a 1:100 dilution to rule out the hook effect.
\end{itemize}

\section*{Normal Results}
\begin{itemize}
  \item \textbf{Prolactin (Adult, non-pregnant female):} 3--25 ng/mL (varies by assay; typical range 4--23 ng/mL for Roche, 2--29 ng/mL for Siemens)
  \item \textbf{Prolactin (Adult, male):} 2--15 ng/mL (slightly lower than female; typical range 4--15 ng/mL)
  \item \textbf{Prolactin (Pregnancy, third trimester):} 50--300 ng/mL (physiologic)
  \item \textbf{Prolactin (Lactation, breastfeeding):} 50--200 ng/mL (basal), with surges to 200--500 ng/mL during suckling
  \item \textbf{Prolactin (Postmenopausal female):} 2--20 ng/mL (slightly lower than pre-menopausal because of lower estrogen)
  \item \textbf{Prolactin (Children, pre-pubertal):} <10 ng/mL
  \item \textbf{Post-PEG precipitation recovery (macroprolactin screen):} >60\% recovery (true monomeric hyperprolactinemia); 40--60\% (indeterminate); <40\% (macroprolactin predominant)
  \item \textbf{Note:} Prolactin reference intervals vary by assay manufacturer and laboratory. Always use the laboratory-specific reference range. Mild prolactin elevation <50 ng/mL in an asymptomatic patient is most commonly due to stress, drugs, macroprolactin, or other secondary causes; a prolactinoma is unlikely if prolactin <50 ng/mL.
\end{itemize}

\section*{Follow-up Tests}
\begin{itemize}
  \item \textbf{Repeat prolactin (rested, fasting, morning sample):} Essential first step for any elevated prolactin. Ideally drawn through an indwelling catheter after 30 minutes of rest to eliminate venipuncture stress. Normalizes in >50\% of mild elevations.
  \item \textbf{Macroprolactin screening by PEG precipitation:} All patients with confirmed hyperprolactinemia should be screened. Post-PEG recovery <40--50\% confirms macroprolactinemia. If macroprolactinemia is confirmed and patient is asymptomatic, no further investigation is needed.
  \item \textbf{TSH and free T4:} To exclude primary hypothyroidism as a cause. Elevated TSH, low FT4, and elevated prolactin that normalizes with levothyroxine therapy confirm the diagnosis. Pituitary imaging is not needed in this scenario.
  \item \textbf{Serum creatinine and eGFR:} To exclude chronic kidney disease as a cause of hyperprolactinemia due to decreased renal clearance.
  \item \textbf{Pregnancy test (beta-hCG):} To exclude pregnancy in any woman of childbearing age with amenorrhea and elevated prolactin.
  \item \textbf{Pituitary MRI with gadolinium contrast:} First-line imaging when true monomeric hyperprolactinemia is confirmed and secondary causes (drugs, hypothyroidism, CKD, pregnancy, macroprolactin) are excluded. Thin-section (1--3 mm) coronal and sagittal T1-weighted images before and after contrast. Detects prolactinoma (micro- or macro-) and other sellar masses. Prolactinoma appears as a hypointense (non-enhancing) area relative to the normally enhancing pituitary on T1 post-contrast images.
  \item \textbf{Diluted prolactin (1:100 dilution):} To exclude the hook effect in patients with macroadenomas (>3 cm) and normal or mildly elevated prolactin --- an absolute necessity before concluding that a large sellar mass is a non-functioning adenoma.
  \item \textbf{Visual field testing (Goldmann or Humphrey perimetry):} In patients with macroprolactinomas abutting or compressing the optic chiasm, to detect early bitemporal hemianopsia before it is clinically apparent.
  \item \textbf{Remaining pituitary function testing:} In patients with large sellar masses, evaluate for hypopituitarism: 8 AM cortisol, ACTH, TSH, free T4, LH, FSH, testosterone (male) or estradiol (female, unless postmenopausal), IGF-1.
  \item \textbf{Testosterone (total and free, morning, male):} Assesses the impact of hyperprolactinemia on the gonadal axis. Hyperprolactinemia suppresses GnRH and LH, leading to secondary hypogonadism (low testosterone). Testosterone normalizes with dopamine agonist therapy.
\end{itemize}

\section*{References}
\bibliographystyle{unsrtnat}
\bibliography{references}

\clearpage
\section*{Regulatory Status and Authorization Date}
Regulatory status applies to the specific assay, instrument, manufacturer, and intended use rather than to the test name alone. No single approval date can be assigned to this test category. For a commercial assay, verify the current product in the relevant regulator's database: the U.S. Food and Drug Administration (FDA) clearance, approval, or authorization; India's Central Drugs Standard Control Organisation (CDSCO) licence or permission; the European Union In Vitro Diagnostic Regulation (EU IVDR) CE certificate issued through the applicable conformity-assessment route; or the United Kingdom Medicines and Healthcare products Regulatory Agency (MHRA) registration and UKCA/CE status. Laboratory-developed tests may not have an individual FDA, CDSCO, EU IVDR, or MHRA approval date.
\clearpage
